\chapter{Аналитическая часть}
% в аналитическом разделе указать краткое описание используемых БЯМов;

Конечный автомат -- модель с конечным набором состояний, переходами между ними по входным символам, начальным состоянием и  принимающими состояниями. Детерминированный конечный автомат (ДКА) для каждого состояния и символа имеет ровно один переход, что позволяет эффективно обрабатывать последовательности.

Регулярное выражение -- формальная нотация для описания регулярных языков, эквивалентная конечным автоматам. Оно строится из базовых операций: конкатенация, альтернатива (|), звезда Клини (*), используя символы алфавита~\cite{lit2}.

\section{БЯМ}

Для получения реализаций по заданию использовались следующие языковые модели~\cite{lit1}:
\begin{itemize}
	\item qwen3-coder -- модель от Alibaba, оптимизированная для кодинга и агентных задач с длинным контекстом, обеспечивает высокую производительность в генерации и анализе кода;
	\item devstral-2 -- Agentic LLM для инженерии ПО, превосходит в работе с инструментами, исследовании кодовых баз и редактировании нескольких файлов;
	\item cogito-2.1 -- модель для когнитивных задач и RAG-систем) с фокусом на обработку знаний;
	\item gemma3 -- компактная модель от Google, эффективна для общих задач ИИ, включая кодинг и рассуждения.
\end{itemize}

\section{Задание по варианту}

Проверка на наличие ошибки использования термина <<оглавление>> вместо <<содержание>>.

\section*{Вывод}
%<что сделали в аналите, кратко>

В данном разделе мной были выбраны БЯМ для решения задания по варианту.

\clearpage
